\worksheettitle
  {Разрядные единицы}
  {\modulemark{M02\(\star\)} Усложнение}

\studentnameblock

\begin{notebox}[Новая ситуация]
  В условии может быть больше девяти одинаковых разрядных единиц. Например,
  \(34\) сотни \textemdash{} это \(34\cdot100=3\,400\). Сначала найдите значение каждой
  части, затем сложите значения.
\end{notebox}

\section{Переводим каждую часть}

\begin{problembox}[Разберите пример]
  Какое число получится из \(6\) десятков тысяч, \(34\) сотен и \(8\) единиц?
  \begin{enumerate}[label=\textbf{\arabic*.}]
    \item \(6\) десятков тысяч: \(6\cdot10\,000=\answerblank[25mm].\)
    \item \(34\) сотни: \(34\cdot100=\answerblank[25mm].\)
    \item Сложите значения:
      \[
        \answerblank[28mm]+\answerblank[28mm]+8=\answerblank[32mm].
      \]
  \end{enumerate}
\end{problembox}

\begin{warningbox}[Не приписывайте цифры]
  Десяток тысяч и сотня \textemdash{} разные разрядные единицы. Нельзя получить ответ,
  просто приписав запись \(34\) к записи \(6\).
\end{warningbox}

\begin{practicebox}[Выполните по образцу]
  \begin{enumerate}[label=\textbf{\arabic*.}]
    \item \(4\) десятка тысяч, \(27\) сотен и \(5\) единиц;
    \item \(3\) сотни тысяч, \(25\) тысяч и \(14\) десятков;
    \item \(8\) десятков тысяч, \(6\) тысяч и \(32\) единицы.
  \end{enumerate}
  \writinglines{5}
\end{practicebox}

\section{Сравниваем разные представления}

\begin{fillbox}[Равны ли числа?]
  Найдите значение каждой записи и поставьте знак \(=\) или \(\ne\).
  \begin{enumerate}[label=\textbf{\arabic*.}]
    \item \(6\) десятков тысяч и \(34\) сотни
      \answerblank[12mm]
      \(63\) тысячи и \(4\) сотни;
    \item \(5\) сотен тысяч и \(28\) сотен
      \answerblank[12mm]
      \(502\) тысячи и \(8\) десятков;
    \item \(7\) десятков тысяч и \(16\) десятков
      \answerblank[12mm]
      \(70\) тысяч и \(160\) единиц.
  \end{enumerate}
\end{fillbox}

\begin{practicebox}[Обратные задачи]
  \begin{enumerate}[label=\textbf{\arabic*.}]
    \item Сколько сотен содержится в числе \(4\,700\)? \answerblank[20mm].
    \item Сколько тысяч содержится в числе \(62\,000\)? \answerblank[20mm].
    \item Сколько десятков содержится в числе \(3\,540\)? \answerblank[20mm].
    \item Представьте \(28\,600\) как сумму десятков тысяч, сотен и единиц.
      \writinglines{1}
  \end{enumerate}
\end{practicebox}

\begin{warningbox}[Проверьте чужое решение]
  Ученик рассуждал так: «\(4\) десятка тысяч и \(27\) сотен \textemdash{} это
  \(427\), потому что я приписал \(27\) к \(4\)».

  Объясните ошибку и запишите верное число.
  \writinglines{1}
\end{warningbox}

\begin{practicebox}[Составьте собственный пример]
  Придумайте условие, в котором названо больше девяти сотен или десятков, а в
  полученном числе появляются нули. Запишите вычисление и ответ.
  \writinglines{2}
\end{practicebox}

\begin{checkpointbox}[Обобщение]
  Почему перед записью числа нужно переводить все названные части в их
  числовые значения?
  \writinglines{1}
\end{checkpointbox}
